\section{Data availability}
\label{sec:availability}

Every artifact named anywhere in this paper is listed here with its path, so
that no repository location has to be carried in the narrative.

\subsection{Preregistration}
The study was preregistered and frozen before any outcome was accessed, under the
identifier \texttt{P1.hidden-formulation.v1}.  A reader checking whether a
reported analysis was planned or chosen after the fact should compare against
that freeze; it is the reason every number reported here can be traced to a
decision made in advance of seeing it.

The machine-readable freeze is \idt{protocol/PROTOCOL_V1.json} and the case
requirements are \idt{protocol/HIDDEN_SHIFT_CASE_SCHEMA_V1.json}. The freeze
records that no outcome had been accessed when it was cut. The amendment that
replaced human adjudication with a blinded model panel is recorded as version
1.1 of that freeze; it carries an adjudication tier that the record type refuses
to report as human adjudication. The adjudication standard itself is
\idt{protocol/ADJUDICATION_RUBRIC_V1.md}, the execution identities are in
\idt{protocol/EXECUTION_MANIFEST_V1.md}, and the prospective power calculation
is \idt{protocol/PROSPECTIVE_POWER_V1.md}.

The archives for Sections~\ref{sec:necessity} and~\ref{sec:revision} label the
following records as their preregistrations:
\idt{research/revival/p1/protocol/P1.epistemic-mutation-necessity.v2.2.4.json}
for the mutation-necessity experiment, together with its superseded earlier
versions in the same directory, and
\idt{research/claim_expansion/p1/P1_X_PROTOCOL_V1.md} with
\idt{research/claim_expansion/p1/v2/P1_X_V2_REPLICATION_PROTOCOL.md} for the
revision-responsibility battery.
These labels and bound bytes identify the registered designs. They do not, by
themselves, independently confirm the original order of protocol and outcome
access.

\subsection{Frozen suite and generated quantities}
The case suite is \idt{protocol/cases/}, split into a pilot and a test set and
bound by content hash: the pilot fingerprint is
\texttt{\PilotSuiteFingerprintShort} and the test fingerprint is
\texttt{\TestSuiteFingerprintShort}. Case generation and validation are
\idt{protocol/cases/generate_cases.py} and
\idt{protocol/cases/validate_cases.py}. Every count, tier and margin quoted in
the text is a macro generated into \idt{manuscript/generated/suite_facts.tex}
by \idt{scripts/render_suite_facts.py}, so prose cannot drift from the frozen
suite. The precision tier the frozen power record assigns to the achieved
half-width is registered there under the label \texttt{\AchievedTier}, with
\RequiredNForHOne{} cases required for the primary hypothesis and
\RequiredNForHTwo{} for the restraint hypothesis. The tracked
\idt{results/INTEGRITY_NOTE.md} preserves the failed historical attempt at
commit \texttt{85d32db0}: 2,794 scored records and 86
\textsc{Cannot-check} records, all in \texttt{orion\_live\_provider} (78
authentication, four rate-limit and four internal-server errors), with 23 of
48 provider cases lacking a scored repeat. That 86-record count describes a
historical execution failure, not current missingness, and no per-cell rerun
manifest is available in a clean checkout. The current tracked raw archive is
authoritative and contains 2,880 of 2,880 \texttt{OK} run records and 2,880 of
2,880 \texttt{SCORED} score records, with zero current
\textsc{Cannot-check} records.

\subsection{The hidden-formulation study}
The scored archive is \idt{results/raw/test_scored.jsonl} with its run records
in \idt{results/raw/test_runs.jsonl}. Table~\ref{tab:P1-T2} projects
\idt{results/P1-T2_baseline_ablation_results.json} and its Markdown companion,
which also carry the full system identifiers, per-family results and mechanistic
metrics; Table~\ref{tab:P1-T3} projects
\idt{results/P1-T3_failure_taxonomy.json} and its Markdown companion, which
carry the per-case detail behind the failure-mode roll-up. The integrity note
for these records is \idt{results/INTEGRITY_NOTE.md}.

The mechanical attribution result and its rule-level limitations are
\idt{evidence/MECHANICAL_ARM_LIMITATIONS_V1.md} with the case table
\idt{evidence/MECHANICAL_ARM_CASE_TABLE_V1.json}. The blind solvability audit is
\idt{evidence/MECHANICAL_SOLVABILITY_AUDIT_V1.md} with
\idt{evidence/mechanical_solvability_audit_v1.json}. The falsifier record,
including the subject commit \texttt{8a8a7feed588363f8e2cd820d3399a33b7af3074}
and continuous-integration run \texttt{31933432314} at which the suite passed,
is \idt{evidence/FALSIFIER_V1.md}. The bounded parent-domain replay is
\idt{evidence/PARENT_DOMAIN_REPLAY_V1.md}, and the saturation record for the
suite is \idt{evidence/FINAL_SATURATION_AUDIT.md}.

\subsection{The mutation-necessity experiment}
The primary and replication directories under
\idt{research/revival/p1/confirmatory/v2.2/} each contain compressed public
worlds, protected response matrices, raw results, the result record, an
separate-implementation verification record, provenance, and a \idt{SHA256SUMS}
manifest.
The archived primary execution identity is \idt{PRIMARY_EXECUTION_FREEZE_V3.json};
earlier receipts and both additive amendments are retained beside it. The
replication directory binds its own disjoint world and execution receipts.
\idt{PRIMARY_REPLICATION_CONCORDANCE.json} checks the cross-run conditions
described as predeclared within the archive without pooling the runs. Because
three cited provenance commits are unreachable, the current audit cannot
independently establish that those records preceded outcome access. The separate verifier is
\idt{research/revival/p1/verify_mutation_necessity_independent.py}; it uses only
the standard library and imports neither the candidate policies, the production
scorer and statistics, nor the subject's licensing code. This supplies code-path
separation within the repository, not external scorer custody.

The current negative-disposition audit is
\idt{revival/r1-negative-revival-audit/AUDIT_RECEIPT.json}, with SHA-256
\idt{e03f3c5b545aa82586539f1d794197bef3674250f42eea5868d8f1c79adbc4ab}.
It confirms 2,882 worlds and 40,348 score rows in each run, zero score or
analysis mismatches, matching supported terminals, and zero task-ID overlap.
It also preserves the broad historical result: subject and strongest baseline
each solve 1 of 48 root tasks, for difference zero.

The clean-source cluster replay is recorded in
\idt{revival/r1-negative-revival-audit/lunarc-replay-v1/LUNARC_REPLAY_RECEIPT_V1.json},
with SHA-256
\idt{f049d2c19c508940188f806315d9b467ccb3139da0ecb2b22e96d17921161acc}.
Job 3550083 ran on node \texttt{cn063} and completed with exit \texttt{0:0} in
\texttt{00:03:45}. The compact job output has SHA-256
\idt{5cd24e6e9a64a9f6b8c4d8c9e89cd9be25cf330352b28b0859b614f40ff28087}.
For both runs, the raw results, result record, and independent-verification
record are byte-identical to the archive. The replay used a Git archive of
source commit \texttt{58192a3a05c4fc8e55b45e0082cbd1de2792c250}; reconstructed
large outputs are not committed.

The same receipts preserve a chronology limitation. Historical execution
records cite commits
\texttt{8c4ac19bb30b93cbb81c90b3fa76cf14b3905418},
\texttt{a7f68bf0fdbfd769d80e4475f526feb23df365c8}, and
\texttt{d3c407b7bc007dabf9c99b1455156e861dfa9a92}, none of which is reachable.
Current deterministic replay therefore does not establish the original
prospective order or external authority.

The round-by-round donor dispositions that fixed the parent and ablation
structure are \idt{research/revival/p1/P1_DONOR_ASSIMILATION_LEDGER_V1.json}
with its later rounds in the same directory,
\idt{research/revival/p1/P1_DONOR_ENGULFMENT_V1.json}, and the saturation record
\idt{research/revival/p1/literature/SATURATION_LEDGER.md}.

\subsection{The revision-responsibility battery}
The protected freeze, protected results, and separate-implementation verification
for both executions are under \idt{research/claim_expansion/p1/}: the first execution in
\idt{results/P1_X_PROTECTED_RESULT_V1.json} with its outcome-access receipt, and
the comparator-fair disjoint second execution in
\idt{v2/P1_X_PROTECTED_RESULT_V2.json} with
\idt{v2/P1_X_INDEPENDENT_VERIFICATION_V2.json}. The case schema is
\idt{research/claim_expansion/p1/SCIENTIFIC_REVISION_RESPONSIBILITY_CASE_V1.schema.json}.

\subsection{Bounded meta-discovery studies}
The study calculus and the three decomposition pilots are under
\idt{development/recursive-atom-calculus-closure-v1/} and
\idt{development/jump-atom-pilot-closure-v1/}. The scoped failure-knowledge
study is \idt{development/failure-epistemology-prospective-closure-v1/}, and the
protected zero-error representation-transition study is
\idt{research/extensions/orion-jump-recursive-atoms/zero_error_jump/}.
The recursive-scientific-evolution donor subtraction is bound by
\idt{RSE_SUCCESSOR_BOUNDARY_V1.md} and records that a generic justification
language is sufficient on the four registered state-schema discriminator pairs.
This is scope narrowing, not superiority evidence.

\subsection{Claim ledgers and nearest work}
The mapping from each result-bearing statement in this manuscript to the
archived artifact that supports it is \idt{evidence/CLAIM_LEDGER.md}, with the
immutable record of the earlier revision at \idt{evidence/CLAIM_LEDGER_V1.md}.
The full nearest-work mechanism matrix, from which
Section~\ref{sec:relatedwork} is drawn, is
\idt{evidence/NEAREST_WORK_MATRIX_V2.md} with
\idt{evidence/nearest_work_matrix_v2.json}. The bounded readiness attestation is
\idt{evidence/PEER_REVIEW_READY_BOUNDED_V2.md}.

\subsection{Naturalistic and comparator preflights}
The outcome-blind metadata snapshot used to prepare the R7A successor is
\idt{development/naturalistic-panel-harness-2026-08-23/SOURCE_METADATA_SNAPSHOT_V1.json},
with canonical payload digest
\idt{7dd283acd81e63e84887d82ee7761e89e62a47d31afceca8628e456b76f73e01}.
Its protocol-binding status and exact non-promotion boundary are recorded in
\idt{development/naturalistic-panel-harness-2026-08-23/PROTOCOL_BINDING_STATUS_V1.json}.
These files contain repository/revision/licence metadata only and do not contain
R7 cases, labels, pair assignments, scorer outputs or outcomes.
The provider- and modality-diverse successor frame is
\idt{development/provider-diverse-metadata-census-2026-08-23/SOURCE_CENSUS_V1.json},
with its frozen rules in
\idt{development/provider-diverse-metadata-census-2026-08-23/CENSUS_PROTOCOL_V1.json}
and its live conformance receipt in
\idt{development/provider-diverse-metadata-census-2026-08-23/VERIFICATION_RECEIPT_V1.json}.
It likewise contains no case or
outcome authority; historical-byte, content-rights and eligibility terminals
remain explicit.

A prospectively frozen public correction/retraction source-feasibility probe is
under
\idt{development/p1-naturalistic-transport-successor-2026-08-23/}. The
protocol, count-only Europe PMC receipt, Crossref rights amendment, source
audit, and checksum manifest are preserved there. The pinned Retraction Watch
CSV itself is not redistributed in this repository; its receipt records 71,944
rows and the byte digest
\idt{ceaab201d728dfcf9929ec1e229acd2ad88c650c847ec922ba9ffe831e366abb}.
No Europe PMC result row, relation, case text, pair assignment, typed-action
gold item, system output, or protected score was opened or created by that
probe. The official CC0 binding applies to the database metadata only, not to
linked article, notice, attachment, abstract, supplement, or blog content.

The outcome-blind rights/relation successor is preserved under
\idt{development/p1-rw-rights-relation-census-v2-2026-08-23/}. Its exact
Retraction Watch relation census, bounded Europe PMC rights reducer,
source-native aggregate strata, amendments, negative ledger, validation
receipt, and 24-entry checksum manifest are retained there. The raw
65,984,968-byte transport is not redistributed; its SHA-256 is
\idt{ceaab201d728dfcf9929ec1e229acd2ad88c650c847ec922ba9ffe831e366abb}.
The deterministic reducer reports 49,878 admitted relations in 42,924
families, including 12,038 exact-both-endpoint-rights relations in 11,602
families. It accessed no case text, scientific terminal, model output, or
protected score. Missing provider matches remain \textsc{Cannot-check}; they
are not incompatible-rights or source-absence findings.

The structured-action semantics successor is
\idt{development/p1-structured-action-semantics-v3-2026-08-23/}. Its frozen
protocol and parser, immutable raw-interface failure, prospectively frozen
normalization amendment, normalized provider census, twelve-group owner-field
feasibility record, rights/provenance receipt, construct-validity failure atlas,
negative ledger, result, validation receipt, and checksum manifest are retained
there. The result digest is
\idt{7178052320f85ea112e4b0b78da558fc5a68f7dc3b465a4f7afc9e1915860f2b}.
The packet retains aggregate token frequencies only: no identifiers, raw
provider responses, case text, scientific gold, model output, or protected
score. Its 52 scientific-artifact assertions pass. The initial one-of-11,991
raw interface match is retained as a serialization failure rather than
overwritten; only the separately frozen display-token normalization supports
the 11,988-of-11,991 relation-set agreement result.

\subsection{Public-standard owner-algebra feasibility}
The outcome-blind standards packet is
\idt{development/p1-owner-algebra-construct-validity-v4-2026-08-23/}.
Its protocol and parser-only amendment were frozen before the corresponding
captures. The retained artifacts bind eight HTTP-200 documents from four
institutional families and one fail-closed HTTP-404 route; raw source documents
were deleted. The packet-native validator checks 85 scientific assertions
without network access, repository CI, case text, outcomes, or protected
scores. Its result records 9/12 structural analogues, 0/12 named-custodian or
delegation groups, 0/12 sufficient groups, and zero scientific-action gold.
The validator and 18-entry manifest establish bounded internal consistency and
byte integrity, not independent semantic review or owner authority.

\subsection{Outcome-blind source-native adapter audit}
The current verifier and its tracked audit receipt are under
\idt{revival/r1-negative-revival-audit/}. They enumerate all 117,649 maps from
repository-held sources, reject 116,929, and reproduce the exact 720-row
survivor registry. Zero maps are fully certified. All 720 survivors are
re-auditable in-repository, but 0/720 are resolvable there. The historical
47-entry external handoff and its checksum receipt remain unavailable in a
clean checkout; they are not required for the current enumeration and provide
no current independent custody.

Every survivor remains blocked on external owner or formally delegated
custodian semantics, target-corpus rights, host authority, a complete signed
seven-target algebra, authorized delivery, and independent semantic review.
V13 remains at 0/7 signed outputs and 0/4 closed authority acts. Zero certified
maps does not mean zero possible maps. The full-withdrawal counterexample is
conditional on later owner-ratification of proposed closed-world target
denotations. This is an external-authority uncertainty, not a setup failure or
an impossibility result.

The external comparator identity and interface audit is
\idt{development/comparator-binding-harness-2026-08-23/COMPARATOR_BINDING_LEDGER_V1.json},
with its scientific interpretation in
\idt{development/comparator-binding-harness-2026-08-23/SCIENTIFIC_REPORT.md}.
Its cross-paper terminal is \idt{CANNOT_CHECK_FULL_COMPARATOR_FREEZE}; bounded
installation or entrypoint probes in that harness grant no P1 performance or
external-authority claim.

\subsection{Owner/custody execution packets V11--V13}
The historical-governance audit is preserved under
\idt{development/p1-owner-custody-positive-successor-v11-2026-08-24/}.
Its \idt{SHA256SUMS} manifest digest is
\idt{8ebd874f0b0a7f3a3a7877c698dee1371352ebc4ce5dc3bb1a1bd2a2d0629721}.
The public issue-chronology, delegation-field and exact-commit route audit is
preserved under
\idt{development/p1-owner-custody-positive-successor-v12-2026-08-24/}.
Its manifest digest is
\idt{7a136154d7fbc3c52d51ea2dade7829446842c21381f3c9c912189d47dc29407}.
The content-addressed external-execution packet is preserved under
\idt{development/p1-owner-custody-positive-successor-v13-2026-08-24/}.
Its manifest digest is
\idt{c2ff7a87356559bfce4ea4033aae807a6a5c4ce03b8b43c29782b58c29d33667}.
These three manifests verify the retained packet bytes only; they do not supply
external identity, delegation, signatures, rights, custody, review, map
certification or scientific outcomes.
\begin{center}
\begin{minipage}{0.96\linewidth}
\textbf{Exact V13 terminal:}\par
\small\raggedright
\idt{P1_V13_CONTENT_ADDRESSED_EXTERNAL_EXECUTION_PACKET_COMPLETE__ZERO_OF_SEVEN_EXTERNAL_SIGNED_OUTPUTS_RECEIVED__ZERO_OF_FOUR_AUTHORITY_ACTS_CLOSED__720_MAPS_UNCHANGED}
\end{minipage}
\end{center}

\section{Code availability}

\subsection{Current review source and historical package identity}
The current review source is not yet bound to a public archive URL, archive
DOI, immutable submission commit, or verified repository-level licence. Editors
and reviewers must therefore receive the exact current snapshot before external
review. The PDF under \idt{journal_package/} remains a content-bound record of
an earlier inspected manuscript, not a rendering of this revision:
\idt{journal_package/RENDER_CLOSURE_STATE.json} marks that package
\texttt{SUPERSEDED}. It must not be relabelled or submitted as the current PDF;
the current source requires a fresh immutable release, build, and visual audit.

\subsection{Repository access and reproduction}
The implementation repository holds the source code, protocol definitions, case
suite, and repository-held evidence artifacts used by this revision. Persistent
DOI assignment remains pending publication. Result-bearing artifacts are identified by their repository paths
and content bindings rather than by an unverified archival claim.

The exact decompression, independent-rescoring, byte-comparison, and
figure-regeneration commands are given in \idt{REPRODUCE.md}. The command
\begin{verbatim}
make paper01-results
\end{verbatim}
regenerates the publication tables from archived raw records only. It does not
execute systems under test or create missing external evidence. On a checkout
without an admissible archived study run the underlying module reports the
result as undetermined; malformed or unbound archives are refused rather than
converted into publishable numbers.

A fresh code and test environment requires Python 3.11 or later, an editable
install with the development extras from the repository root, and
model-provider access only for commands that actually execute a live-provider
arm; dependencies are declared in \idt{pyproject.toml}.

\subsection{Runtime and cost}
The credential-free mutation-necessity campaign and its independent
verification each complete in approximately one minute on the current CPU
environment. They make no network call and incur no model-provider cost.
Algorithmic effort is fixed at 2,882 worlds, nine runnable arms, five ablations,
and 10,000 registered paired-bootstrap resamples. The provider cost of the
earlier live-provider arm was not recorded under a bound environment and
therefore remains undetermined; it is not imputed from the mechanical
successor.

\subsection{Licensing status}
This revision does not contain a repository-level licence file or a project
licence declaration in \idt{pyproject.toml}. Redistribution terms for code, case
definitions, and protocol artifacts therefore remain undetermined from the
repository state and must be resolved explicitly before archival or
publication. Model-provider terms remain separately governed by the relevant
providers.

\subsection{Artifact fingerprinting}
Result-bearing execution identities are bound through the protocol and archived
records. Relevant coordinates include dataset and suite revisions and
fingerprints; baseline configuration hashes; evaluator identity; frozen split
hashes; and the exact subject revision and provider/model revisions for a live
run. Reproduction of an empirical result requires the exact bound subject and
admissible archived records; prose or a protocol declaration alone is not
evidence that the run occurred.
